% ===================================================================
%  TAULA N.º 1 — L'Escola. — L'Institut. — L'Aula.
%
%  Ceci n'est PAS une transcription : c'est une traduction, faite en
%  2026, en catalan d'aujourd'hui. D'ou trois differences avec
%  texto/io et texto/fr, et trois seulement :
%
%   * pas de \nl ni de \cc — la traduction ne suit aucune ligne
%     imprimee, et les lignes de ce fichier ne sont que du confort ;
%   * pas de \begin{VUpage} — elle n'a ni feuillet ni folio, et la
%     page de lecture ne lui en affiche pas ;
%   * le gras se pose sur le TERME seul, ponctuation dehors.
%
%  Tout le reste est identique, et doit l'etre : les cles « %%K » sont
%  celles de texto/io, macro pour macro pour l'apparat, et les renvois
%  \textsuperscript{(n)} visent les MEMES objets de la planche — ce
%  sont eux qui ouvrent les gros plans.
%
%  LA SIXIEME DES DIX-SEPT LANGUES IDISTES. La Catalogne a eu ses
%  groupes idistes des les annees 1910. La colonne se traduit sur
%  l'ido, et se controle sur texto/es, qui est la plus proche des
%  colonnes deja faites : le catalan et le castillan partagent l'ordre
%  des mots et presque tout le lexique de ce livret.
%
%  MAIS ILS NE LE PARTAGENT PAS TOUT ENTIER, et c'est la qu'il faut
%  se mefier. « pizarra » se dit « pissarra », mais « borrador » se
%  dit « esborrador » et « tiza » se dit « guix » ; « plumier » se dit
%  « plomer », « atril » se dit « faristol », « percha » se dit
%  « penjador ». La ressemblance est une aide, pas une regle.
%
%  L'ARTICLE SE CONTRACTE ET S'APOSTROPHE : « de el » donne « del »,
%  « a el » donne « al », et « la aula » donne « l'aula ». Le mot qui
%  porte le renvoi n'en est pas touche — l'apostrophe se fait devant,
%  jamais derriere — et l'ordre des renvois n'en souffre pas.
%
%  LES NOMS DE PERSONNE SONT CATALANS, comme dans les colonnes
%  precedentes : Ioannes est Joan, Iakobus est Jaume, Karolus est
%  Carles. Seuls le chien Medor et l'ane Aliboron gardent leur nom.
%  Les noms latins de la note (*) ne bougent pas : c'est d'eux que la
%  note parle.
% ===================================================================

%%K t01-apar-0 apar
\VUsaut{2.0mm}
\VUcentre{12.6pt}{120}{FULLET EXPLICATIU}
\VUsaut{3.2mm}
\VUcentre{8.4pt}{160}{DE}
\VUsaut{2.6mm}
\VUtitre{15.6pt}{0}{les Taules Auxiliars Delmas}
\VUsaut{3.4mm}
\VUornamento{9pt}
\VUsaut{5.0mm}
\VUcentre{13.2pt}{90}{PRIMERA SÈRIE}
\VUsaut{3.0mm}
\VUfilet{16mm}
\VUsaut{5.6mm}

%%K t01-apar-1 apar
\VUtitre{15.0pt}{60}{TAULA N.º 1}
\VUsaut{1.4mm}
\VUtitre{11.4pt}{0}{L'Escola. --- L'Institut. --- L'Aula.}
\VUsaut{2.2mm}
\VUfilet{20mm}
\VUsaut{6.4mm}

%%K t01-tit-1 sub
\VUpk{10.2pt}{40}{Descripció general.}
\VUsaut{3.0mm}

%%K t01-01-1 p
1. --- Aquesta taula representa una \VUgras{aula} d'\VUgras{institut}. En
aquesta aula hi ha vuit \VUgras{taules} \textsuperscript{(18)} i vuit
\VUgras{bancs} \textsuperscript{(32)}. En cada banc s'asseuen tres alumnes. El
\VUgras{professor} \textsuperscript{(7)} està assegut en una \VUgras{cadira}
\textsuperscript{(8)} a la seva \VUgras{càtedra} \textsuperscript{(6)}.

%%K t01-02-1 p
2. --- A l'esquerra de la càtedra hi ha un \VUgras{armari} amb portes de vidre
\textsuperscript{(15)}. A la dreta hi ha la \VUgras{pissarra}
\textsuperscript{(1)} amb l'\VUgras{esborrador} \textsuperscript{(2)} i el
\VUgras{guix} \textsuperscript{(3)}.

%%K t01-02-2 p
Damunt de la càtedra pengen unes \VUgras{làmines murals}
\textsuperscript{(9, 11, 12)} i un \VUgras{mapa} \textsuperscript{(10)}.

%%K t01-03-1 p
3. --- No tots els alumnes són al seu \VUgras{lloc} \textsuperscript{(17)}. Joan
\textsuperscript{(4)} és dret vora la pissarra; sosté l'esborrador i esborra les
\VUgras{figures de geometria}.

%%K t01-03-2 p
Pere és a prop de la càtedra; recull els \VUgras{exercicis escrits}
\textsuperscript{(5)} i els porta al professor.

%%K t01-04-1 p
4. --- \VUgras{Aquest} professor és el professor de \VUgras{llengües modernes}.
Ensenya els alumnes a parlar, llegir i escriure llengües estrangeres
\VUgras{(alemany, anglès, espanyol, italià, rus} o \VUgras{francès)}.

%%K t01-05-1 p
5. --- L'aula és molt àmplia i molt clara. Al fons hi ha una \VUgras{finestra}
ampla amb dos \VUgras{muntants} \textsuperscript{(78)} i una \VUgras{porta
envidrada} per ventilar-la i il·luminar-la.

%%K t01-05-2 p
Aquesta aula no és freda. Al mig, per escalfar-la, hi ha una \VUgras{estufa}
\textsuperscript{(46)} molt bona amb el seu \VUgras{tub} \textsuperscript{(45)}.

%%K t01-05-3 p
A l'envà hi ha fixats uns \VUgras{penjadors} \textsuperscript{(58)}; d'ells
pengen les gorres i els abrics dels alumnes.

%%K t01-tit-2 sub
\VUsaut{5.2mm}
\VUpk{10.2pt}{40}{El pati d'esbarjo.}
\VUsaut{3.6mm}

%%K t01-06-1 p
6. --- El \VUgras{rellotge} \textsuperscript{(63)} marca les nou i cinc.

%%K t01-06-2 p
L'\VUgras{esbarjo} \textsuperscript{(76)} acaba d'acabar i la classe torna a
començar. L'esbarjo es fa al \VUgras{pati} \textsuperscript{(75)}. Veiem uns
quants alumnes que juguen. Un \VUgras{passant} \textsuperscript{(74)} els
vigila.

%%K t01-07-1 p
7. --- Al pati veiem a més altres persones: l'\VUgras{inspector}
\textsuperscript{(73)}, el \VUgras{director} \textsuperscript{(71)} i el
\VUgras{cap d'estudis} \textsuperscript{(72)}.

%%K t01-07-2 p
L'inspector inspecciona les escoles, el director dirigeix l'institut, el cap
d'estudis vigila els estudis.

%%K t01-07-3 p
El quart personatge és el \VUgras{conserge} \textsuperscript{(77)}. Porta sota el
braç un registre per anotar-hi els absents.

%%K t01-08-1 p
8. --- Al fons del pati hi ha els \VUgras{aparells de gimnàstica}
\textsuperscript{(70)} i el taller de \VUgras{treballs manuals}
\textsuperscript{(69)}.

%%K t01-08-2 p
A la dreta de la taula comencem a veure les \VUgras{aules} de \VUgras{música} i
de \VUgras{dibuix}.

%%K t01-noto-1 noto
\VUnotes{143.2mm}{%
(*) Per als \textit{noms de pila} s'aconsella transcriure'ls també en la seva
forma llatina. És l'única manera d'evitar variants sense compte. D'altra banda,
com triar entre \textit{Giovanni, Jean,} \textit{Johann, Jan, Hans, John, Ivan,}
etc.? Caldrà crear un diccionari especial de noms de pila? Prenguem del llatí el
seu nominatiu i diguem: \VUgras{Ioannes, Ioanna; Iulius, Iulia; Iakobus;
Andreas; Lukas.} (Gramàtica completa).%
}

%%K t01-tit-3 sub
\VUpk{10.2pt}{40}{Descripció detallada.}
\VUsaut{2.4mm}

%%K t01-tit-4 sub
\VUpk{10.2pt}{40}{Els bons alumnes.}
\VUsaut{3.0mm}

%%K t01-09-1 p
9. --- Els alumnes del primer banc a la dreta de la càtedra són \VUgras{Enric}
(*), \VUgras{Esteve} i \VUgras{Carles}.

%%K t01-09-2 p
Enric és molt atent i treballador; escolta el professor. Sobre la seva taula té
un \VUgras{quadern} \textsuperscript{(28)}, un \VUgras{llibre}
\textsuperscript{(29)}, un \VUgras{diccionari} \textsuperscript{(30)} i un
\VUgras{plomer} \textsuperscript{(31)}. El seu llibre té moltes
\VUgras{pàgines}; cada capítol conté diversos \VUgras{paràgrafs} i cada
\VUgras{línia} moltes \VUgras{paraules}.

%%K t01-09-4 p
El seu veí Esteve \textsuperscript{(24)} és aplicat. Apunta al seu quadern la
lliçó per a la pròxima vegada. Té bona \VUgras{lletra}. El seu llibre és tancat.
Només veiem la \VUgras{coberta} \textsuperscript{(27)}. Al seu costat hi ha el
seu \VUgras{compàs} \textsuperscript{(26)}.

%%K t01-09-5 p
Carles \textsuperscript{(23)} sosté el seu llibre a la mà; l'obre per repassar la
lliçó. Com tots els alumnes, té al davant un \VUgras{tinter}
\textsuperscript{(25)}. És obedient.

%%K t01-10-1 p
10. --- A la taula del costat hi ha \VUgras{Lluís, Antoni} i \VUgras{Pau}. Lluís
recita la seva \VUgras{lliçó} \textsuperscript{(22)} i respon a les
\VUgras{preguntes} del professor. Antoni \textsuperscript{(21)} està molt
distret; de vegades el professor fins i tot el castiga. Pau
\textsuperscript{(20)} és bon company, amable i servicial. És molt atent.

%%K t01-tit-5 sub
\VUsaut{4.2mm}
\VUpk{10.2pt}{40}{Els mals alumnes.}
\VUsaut{3.2mm}

%%K t01-11-1 p
11. --- Darrere d'Enric hi ha \VUgras{Jordi} \textsuperscript{(38)}. No és mal
noi, però és \VUgras{xerraire}, distret i inquiet. La seva \VUgras{lleugeresa} i
la seva \VUgras{desobediència} causen \VUgras{desordre} durant la classe, i sovint
mereix una severa \VUgras{advertència}. El seu \VUgras{quadern d'esborrany}
\textsuperscript{(35)} és brut, l'omple de \VUgras{taques de tinta} amb la seva
\VUgras{ploma} \textsuperscript{(34)}, i per això sempre fa servir la seva
\VUgras{goma d'esborrar} \textsuperscript{(36)}; la seva \VUgras{cartera}
\textsuperscript{(33)} és trencada, perquè és molt descurat. Per què porta el seu
\VUgras{portallapis} \textsuperscript{(37)} a la classe de llengües modernes?

%%K t01-11-2 p
El seu veí Edmon \textsuperscript{(39)} no l'estima. És cert que és una mica
peresós, però és callat i s'està quiet. No xerra; només que, en comptes
d'atendre, s'estima més llegir els \VUgras{contes} del seu \VUgras{llibre de
lectura}.

%%K t01-11-3 p
El tercer alumne d'aquest banc és \VUgras{Lluís} \textsuperscript{(40)}. És
aplicat. Escriu en un \VUgras{full de paper} blanc. Davant seu ha deixat el seu
\VUgras{paper assecant} \textsuperscript{(41)}.

%%K t01-12-1 p
12. --- Els alumnes del segon banc, a l'esquerra del professor, són molt joves.
El primer, el petit \VUgras{Emili}, encara no sap escriure gaire bé; porta el seu
\VUgras{pissarrí} \textsuperscript{(42)}, el seu \VUgras{llapis de pissarra} i la
seva \VUgras{esponja} \textsuperscript{(43)}.

%%K t01-12-2 p
Al seu costat hi ha \VUgras{August} i \VUgras{Bernat} \textsuperscript{(44)}.
Aquest darrer treballa bé, però el seu \VUgras{aspecte} és molt descurat.

%%K t01-13-1 p
13. --- Darrere d'ells hi ha tres \VUgras{mandrosos}. \VUgras{Albert} no s'aprèn
les lliçons. \VUgras{Jaume} \textsuperscript{(47)} dorm en comptes d'escoltar.
La seva \VUgras{peresa} li comporta \VUgras{reprensions} i fins i tot
\VUgras{càstigs}. Finalment, \VUgras{Cristià} \textsuperscript{(48)} és massa
frívol. Es \VUgras{porta} malament i no fa cap esforç per esmenar-se.

%%K t01-14-1 p
14. --- Els seus veïns, en canvi, es porten bé. \VUgras{Ernest}
\textsuperscript{(51)} estima l'ordre i la regularitat; sempre fa servir el seu
\VUgras{regle} \textsuperscript{(50)} i el seu \VUgras{llapis}
\textsuperscript{(49)}. És \VUgras{extern}.

%%K t01-14-2 p
\VUgras{Francesc} \textsuperscript{(52)} és \VUgras{intern}; porta l'\VUgras{uniforme}
de l'institut; hi menja i hi dorm. És un bon \VUgras{company}.

%%K t01-14-3 p
Maurici esmola el seu \VUgras{llapis} amb la seva \VUgras{navalla}
\textsuperscript{(56)}; és molt curós.

%%K t01-14-4 p
Porta tots els seus llibres, la seva \VUgras{gramàtica} \textsuperscript{(55)},
el seu \VUgras{quadern de notes} \textsuperscript{(54)} i fins i tot una
\VUgras{carpeta d'escriptori} \textsuperscript{(53)}. La seva \VUgras{cartera
d'escola} \textsuperscript{(57)} és a terra darrere d'ell.

%%K t01-14-5 p
Les dues taules del fons de l'aula les ocupen \VUgras{bons alumnes}
\textsuperscript{(19)}.

%%K t01-tit-6 sub
\VUsaut{4.6mm}
\VUpk{10.2pt}{40}{Dibuix i música.}
\VUsaut{3.4mm}

%%K t01-15-1 p
15. --- A l'\VUgras{aula de dibuix} hi ha bonics \VUgras{models} penjats de la
paret i una \VUgras{làmpada} \textsuperscript{(62)} amb \VUgras{pantalla},
penjada del sostre. El \VUgras{professor} \textsuperscript{(60)} és molt
pacient; corregeix els \VUgras{dibuixos} \textsuperscript{(61)} dels alumnes i
els ensenya a dibuixar un \VUgras{bust} \textsuperscript{(59)} del natural.

%%K t01-16-1 p
16. --- L'\VUgras{aula de música} sembla molt interessant. Allà s'aprèn a llegir
la \VUgras{música}, a solfejar les \VUgras{notes de l'escala}
\textsuperscript{(67)} escrites al \VUgras{pentagrama} (do, re, mi, fa, sol, la,
si\,=\,C. D. E. F. G. A. B.), a reconèixer les \VUgras{claus} i a cantar boniques
\VUgras{melodies}. Les partitures es posen al \VUgras{faristol}
\textsuperscript{(65)}. Un dels alumnes \VUgras{toca el piano}
\textsuperscript{(64)} i el \VUgras{professor de música} \textsuperscript{(66)}
\VUgras{porta el compàs} \textsuperscript{(68)}.

%%K t01-17-1 p
17. --- Comencem a veure un racó del \VUgras{taller} de \VUgras{treballs
manuals} \textsuperscript{(69)}, on els nois poden aprendre a raspallar la fusta
i a llimar el ferro. No n'hi ha a tots els instituts.

%%K t01-tit-7 sub
\VUsaut{5.0mm}
\VUpk{10.2pt}{40}{L'aula de ciències.}
\VUsaut{3.6mm}

%%K t01-18-1 p
18. --- El professor de matemàtiques ensenya als alumnes \VUgras{geometria,
aritmètica, càlcul} i el \VUgras{sistema mètric}. Veiem a la taula les principals
figures de geometria: un \VUgras{cercle} \textsuperscript{(a)}, un
\VUgras{quadrat} \textsuperscript{(b)}, un \VUgras{triangle}
\textsuperscript{(c)}, una \VUgras{línia} \textsuperscript{(d)}.

%%K t01-18-2 p
Veiem també una \VUgras{operació}; no és una \VUgras{suma}, ni una
\VUgras{resta}, ni una \VUgras{multiplicació}: és una \VUgras{divisió}
\textsuperscript{(e)}.

%%K t01-19-1 p
19. --- A la làmina del sistema mètric distingim: un \VUgras{metre}
\textsuperscript{(a)}, una \VUgras{balança} \textsuperscript{(b)} i les seves
\VUgras{peses}, un \VUgras{litre} \textsuperscript{(e)}, un \VUgras{decalitre}
\textsuperscript{(f)}, un \VUgras{decilitre} i un \VUgras{metre cúbic}
\textsuperscript{(d)}.

%%K t01-19-2 p
Els alumnes estudien també \VUgras{ciències naturals} \textsuperscript{(11)} ---
zoologia \textsuperscript{(a)}, botànica \textsuperscript{(b)}, geologia
\textsuperscript{(c)} --- i \VUgras{història} \textsuperscript{(12)}.

%%K t01-19-3 p
A l'armari hi ha els instruments de \VUgras{física} \textsuperscript{(15)} i de
\VUgras{química} \textsuperscript{(16)}.

%%K t01-19-4 p
Damunt de l'armari hi ha: una \VUgras{esfera} \textsuperscript{(14)}, un
\VUgras{con} i un \VUgras{globus terraqüi} \textsuperscript{(13)}.

%%K t01-19-5 p
Damunt de les làmines penja un \VUgras{mapa} que representa les cinc parts del
món: \VUgras{Europa} \textsuperscript{(g)}, \VUgras{Àsia}
\textsuperscript{(e)}, \VUgras{Àfrica} \textsuperscript{(f)},
\VUgras{Amèrica} \textsuperscript{(ab)} i \VUgras{Oceania}
\textsuperscript{(h)}, i els dos grans oceans, el \VUgras{Pacífic}
\textsuperscript{(c)} i l'\VUgras{Atlàntic} \textsuperscript{(d)}.
\VUsaut{6.0mm}
\VUfilet{22mm}
